\documentclass[10pt]{article}
\usepackage[T1]{fontenc}
\usepackage[utf8]{inputenc}
\usepackage[french]{babel}
\usepackage{amsmath,amssymb,mathrsfs,mathpazo}
\usepackage{geometry}
\geometry{a4paper,margin=2cm}
\usepackage{graphicx,array,tabularx,tabularray,colortbl,multirow,multicol,booktabs,enumitem,pifont,ulem,fancybox,diagbox,cancel}
\usepackage{pstricks,pstricks-add,pst-plot,pst-tree,pst-node,pst-text,pst-func,pst-math,pst-bspline,pst-3dplot,pst-eucl,pst-solides3d}
\let\pscalculate\undefined
\usepackage{pst-bezier}
\usepackage{tikz,pgfplots,tkz-tab}
\usetikzlibrary{arrows,arrows.meta,patterns,calc,positioning,trees,decorations.markings,shapes.geometric,shapes.misc,shapes.arrows,intersections,decorations.pathmorphing,decorations.pathreplacing,backgrounds,shadows,babel}
\pgfplotsset{compat=1.18}
\usepgfplotslibrary{fillbetween}
\usepackage[np]{numprint}
\usepackage{listings,makecell,pas-tableur,twoopt}
\let\barre\undefined
\usepackage[pstricks]{tablvar}
\definecolor{darkgreen}{RGB}{0,100,0}
\providecommand{\red}{\color{red}}
\providecommand{\blue}{\color{blue}}
\providecommand{\green}{\color{green}}
\providecommand{\black}{\color{black}}
\providecommand{\white}{\color{white}}
\providecommand{\gray}{\color{gray}}
\providecommand{\cyan}{\color{cyan}}
\providecommand{\magenta}{\color{magenta}}
\providecommand{\yellow}{\color{yellow}}
\providecommand{\violet}{\color{violet}}
\providecommand{\orange}{\color{orange}}
\providecommand{\euro}{\texteuro}
\providecommand{\EUR}{\texteuro}
\providecommand{\ieme}{\textsuperscript{e}}
\providecommand{\degres}{\ensuremath{^\circ}}
\providecommand{\degree}{\ensuremath{^\circ}}
\providecommand{\vv}[1]{\overrightarrow{#1}}
\providecommand{\Coord}[1]{#1}
\providecommand{\decosix}{}
\providecommand{\starredbullet}{\ensuremath{\bullet}}
\pagestyle{empty}
\setlength{\parindent}{0pt}
\usepackage{adjustbox}\DeclareUnicodeCharacter{2236}{:}\DeclareUnicodeCharacter{2190}{\ensuremath{\leftarrow}}
\providecommand{\R}{\mathbb{R}}
\providecommand{\N}{\mathbb{N}}
\providecommand{\Z}{\mathbb{Z}}
\providecommand{\Q}{\mathbb{Q}}
\providecommand{\D}{\mathbb{D}}
\providecommand{\C}{\mathbb{C}}
\providecommand{\e}{{\mathrm{e}}}
\providecommand{\vect}[1]{\overrightarrow{#1}}
\providecommand{\vectt}[1]{\overrightarrow{\mathrm{#1}}}
\providecommand{\vvt}[1]{\overrightarrow{\mathrm{#1}}}
\providecommand{\np}[2][]{{#2~#1}}
\providecommand{\vv}[1]{{\overrightarrow{#1}}}
\renewcommand{\d}{{\mathrm{d}}}
\providecommand{\degres}{{^\circ}}
\providecommand{\degre}{{^\circ}}
\providecommand{\blue}{\color{blue}}
\providecommand{\red}{\color{red}}
\providecommand{\enskip}{\quad}
\providecommand{\llbracket}{\lbrack\!\lbrack}
\providecommand{\rrbracket}{\rbrack\!\rbrack}
\providecommand{\euro}{{\text{€}}}
\providecommand{\no}{{\text{n°}}}
\providecommand{\interval}[2]{{[#1;#2]}}
\providecommand{\intervFO}[2]{{[#1;#2[}}
\providecommand{\dotfill}{\ldots}

\providecommand{\ds}{\displaystyle}
\providecommand{\textding}[1]{#1}
\providecommand{\Oij}{\(\mathrm{O};\vec{\imath},\vec{\jmath}\)}
\providecommand{\Oijk}{\(\mathrm{O};\vec{\imath},\vec{\jmath},\vec{k}\)}
\providecommand{\Ouv}{\(\mathrm{O};\vec{u},\vec{v}\)}
\providecommand{\intervFF}[2]{\ensuremath{[#1;#2]}}
\renewcommand{\euro}{\ifmmode\text{\texteuro}\else\texteuro\fi}
\DeclareUnicodeCharacter{00B0}{\ifmmode^{\circ}\else\textdegree\fi}
\DeclareUnicodeCharacter{20AC}{\ifmmode\text{\texteuro}\else\texteuro\fi}
% Énoncés : Éducation nationale. Transcription et corrigés : APMEP, Denis Vergès et contributeurs. Découpage et classement : CoopMaths, Sébastien Lozano, Éric Elter et contributeurs. Adaptation de la présentation : Maths Explorer.
% CC BY-SA : https://coopmaths.fr/www/about/
% https://coopmaths.fr/alea/static/bac/2023/tex/bac_2023_03_sujet1_metropole_1_cor.tex
\begin{document}
\noindent\textit{APMEP / CoopMaths. CC BY-SA. Présentation adaptée par Maths Explorer.}\par\medskip
Un technicien contrôle les machines équipant une grande entreprise. Toutes ces machines sont identiques.

On sait que :

\begin{itemize}
\item[$\bullet~~$]20\,\% des machines sont sous garantie;
\item[$\bullet~~$]0,2\,\% des machines sont à la fois défectueuses et sous garantie;
\item[$\bullet~~$]8,2\,\% des machines sont défectueuses.
\end{itemize}

\smallskip

Le technicien teste une machine au hasard.




On considère les évènements suivants :
\begin{itemize}
\item[$\bullet~~$]$G$: \og la machine est sous garantie \fg{} ;
\item[$\bullet~~$]$D$: \og la machine est défectueuse\fg{} ;
\item[$\bullet~~$]$\overline{G}$ et $\overline{D}$ désignent respectivement les évènements contraires de $G$ et $D$.
\end{itemize}


















\begin{enumerate}
\item La probabilité $p_G(D)$ de l'évènement $D$ sachant que $G$ est réalisé est égale à :
\begin{center}
\begin{tabularx}{\linewidth}{*{4}{X}}
\textbf{a.~~}0,002 &\textbf{b.~~} 0,01 &\textbf{c.~~} 0,024 & \textbf{d.~~} 0,2
\end{tabularx}
\end{center}



\begin{tabular}{@{\hspace*{0.05\linewidth}} || p{0.93\linewidth}}
20\,\% des machines sont sous garantie donc $p(G)=0,2$.

0,2\,\% des machines sont à la fois défectueuses et sous garantie donc $p(G\cap D)=0,002$.

$p_{G}(D) = \dfrac{p(G\cap D)}{p(G)}
= \dfrac{0,002}{0,2}=0,01$
\hfill\textbf{Réponse b.}
\end{tabular}

\item La probabilité $p\left(\overline{G} \cap D\right)$ est égale à :
\begin{center}
\begin{tabularx}{\linewidth}{*{4}{X}}
\textbf{a.~~}0,01 &b. 0,08 &\textbf{c.~~} 0,1  &\textbf{d.~~} 0,21
\end{tabularx}
\end{center}



\begin{tabular}{@{\hspace*{0.05\linewidth}} || p{0.93\linewidth}}
8,2\,\% des machines sont défectueuses donc $p(D)=0,082$.

D'après la formule des probabilités totales:
$p(D)= p(G\cap D) + p\left(\overline{G} \cap D\right)$. Donc:

$p\left(\overline{G} \cap D\right) 
= p(D) - p(G\cap D) = 0,082 - 0,002 = 0,08$
\hfill\textbf{Réponse b.}
\end{tabular}

\item La machine est défectueuse. La probabilité qu'elle soit sous garantie est
environ égale, à $10^{-3}$ près, à : 
\begin{center}
\begin{tabularx}{\linewidth}{*{4}{X}}
\textbf{a.~~}0,01  &\textbf{b.~~} 0,024  &\textbf{c.~~} 0,082 &\textbf{d.~~} 0,1
\end{tabularx}
\end{center}



\begin{tabular}{@{\hspace*{0.05\linewidth}} || p{0.93\linewidth}}
On cherche:
$p_{D}(G)=\dfrac{p(G\cap D)}{p(D)}=\dfrac{0,002}{0,082}\approx 0,024$
\hfill\textbf{Réponse b.}
\end{tabular}

\end{enumerate}

Pour les questions 4 et 5, on choisit au hasard et de façon indépendante $n$ machines de l'entreprise, où $n$ désigne un entier naturel non nul. 

On assimile ce choix à un tirage avec remise, et on désigne par $X$ la variable aléatoire qui associe à chaque lot de $n$ machines le nombre de machines défectueuses dans ce lot.

On admet que $X$ suit la loi binomiale de paramètres $n$ et $p = 0,082$.

\begin{enumerate}[resume]
\item Dans cette question, on prend $n =50$.

La valeur de la probabilité $p(X > 2)$, arrondie au millième, est de : 
\begin{center}
\begin{tabularx}{\linewidth}{*{4}{X}}
\textbf{a.~~}0,136 &\textbf{b.~~} 0,789 &\textbf{c.~~} 0,864 &d. 0,924
\end{tabularx}
\end{center}



\begin{tabular}{@{\hspace*{0.05\linewidth}} || p{0.93\linewidth}}
$p(X>2) = 1-p(X\leqslant 2)\approx 1-0,2114 \approx 0,789$
\hfill\textbf{Réponse b.}
\end{tabular}

\item On considère un entier $n$ pour lequel la probabilité que toutes les machines d'un lot de taille $n$ fonctionnent correctement est supérieure à $0,4$. 

La plus grande valeur possible pour $n$ est égale à :
\begin{center}
\begin{tabularx}{\linewidth}{*{4}{X}}
\textbf{a.~~}5 &\textbf{b.~~} 6 &\textbf{c.~~} 10 &d. 11
\end{tabularx}
\end{center}



\begin{tabular}{@{\hspace*{0.05\linewidth}} || p{0.93\linewidth}}
Si toutes les machines fonctionnent, le nombre de machines défectueuses est 0; donc on cherche le plus grand $n$ tel que $p(X=0)>0,4$.

$p(X=0) = \ds\binom{10}{0}\times 0,082^0 \times (1-0,082)^n = 0,918^n$

$p(X=0)>0,4
\iff
0,918^n > 0,4
\iff
\ln\left (0,918^n\right ) > \ln(0,4)$

$\phantom{p(X=0)>0,4}
\iff
n\times \ln(0,918) > \ln (0,4)
\iff
n < \dfrac{\ln(0,4)}{\ln(0,918)}$

$\dfrac{\ln(0,4)}{\ln(0,918)} \approx 10,71$ donc le plus grand entier $n$ inférieur à $\dfrac{\ln(0,4)}{\ln(0,918)}$ est 10.
\hfill\textbf{Réponse c.}

\end{tabular}

\end{enumerate}

\bigskip

\end{document}
